% Figure-caption inputs for plot geometry.

\begin{figure}[!htbp]
\centering
\includegraphics{figures/fig_budget.pdf}
\caption{Event sensitivity and achieved burden at the two evaluated nominal calibration budgets. \textbf{a}, \textbf{b}, Equal-source participant-macro sensitivity at 30 and 60 minutes. \textbf{c}, \textbf{d}, Equal-source participant-macro unmatched alarms per participant-calendar day at the corresponding horizons. Thresholds for nominal budgets of 0.5 and 1.0 are fitted on calibration participants. Lines join those two evaluated policies and do not represent an additional threshold sweep. Large markers average the five LSTM or TCN fits, with smaller marks showing individual seeds. Each classical arm contributes one fit or deterministic rule. All policies use base eligibility and the same 153 selection participants, 24{,}574 participant-calendar days, and 3{,}616 or 3{,}674 base events at 30 or 60 minutes. A common nominal budget does not imply equal achieved selection burden.}
\label{sfig:budget}
\end{figure}

\begin{figure}[!htbp]
\centering
\includegraphics{figures/fig_frontier.pdf}
\caption{Eligibility trajectories for all six arms on the fixed selection population. Panels pair 30-minute and 60-minute horizons for the current-glucose control (\textbf{a}, \textbf{b}), linear extrapolation (\textbf{c}, \textbf{d}), summary logistic regression (\textbf{e}, \textbf{f}), XGBoost (\textbf{g}, \textbf{h}), LSTM (\textbf{i}, \textbf{j}) and TCN (\textbf{k}, \textbf{l}). Sensitivity averages event-positive participants within source and then weights sources equally. Solid curves use calibration-refitted thresholds at each eligibility cut and dashed curves retain the base calibration threshold. Both use a nominal base budget of 0.5 unmatched alarms per participant-calendar day. Neural curves average five fits, with thin lines showing individual seeds. The classical arms each use one fit or deterministic rule. Events and exposure remain fixed across cuts, including events with no permitted evaluable lead-in. No selection-partition threshold sweep is used.}
\label{sfig:frontier}
\end{figure}

\begin{figure}[!htbp]
\centering
\includegraphics{figures/fig_reliability.pdf}
\caption{Window-level probability diagnostics for all four classification heads on selection participants. \textbf{a}, \textbf{b}, Level 2 hypoglycaemia at 30 minutes. \textbf{c}, \textbf{d}, Level 2 hypoglycaemia at 60 minutes. \textbf{e}, \textbf{f}, Hyperglycaemia at 30 minutes. \textbf{g}, \textbf{h}, Hyperglycaemia at 60 minutes. Left panels show area under the precision--recall curve (AUPRC) with filled markers and area under the receiver operating characteristic curve (AUROC) with open markers. Right panels show Brier score with filled markers and expected calibration error (ECE) with open markers. The latter panels give the actual retained window count and positive-label frequency. Large markers average the five LSTM or TCN fits, with smaller marks showing individual seeds. Each classical arm contributes one fit or deterministic rule. All probabilities use temperatures fitted only on calibration participants. No selection-fitted transformation is applied. Labelled windows may overlap and are not independent event or participant replicates.}
\label{sfig:reliability}
\end{figure}

\begin{figure}[!htbp]
\centering
\includegraphics{figures/fig_seed_variance.pdf}
\caption{Source-resolved model-seed variation under base eligibility. Heatmaps show participant-macro event sensitivity for LSTM at \textbf{a}, 30 minutes and \textbf{b}, 60 minutes, and TCN at \textbf{c}, 30 minutes and \textbf{d}, 60 minutes. Each cell is one actual source--seed--horizon result. The four heatmaps share a zero-to-one scale. Row labels give event-positive participant counts, identical at both horizons, and columns identify the five training seeds. Thresholds are fitted on calibration data at a nominal 0.5 unmatched alarms per participant-calendar day. \textbf{e}, \textbf{f}, Maximum-minus-minimum sensitivity across five seeds for LSTM and TCN, with filled circles for 30 minutes and open squares for 60 minutes. All fits share the same data and participant split. These ranges describe optimisation variation rather than participant sampling uncertainty.}
\label{sfig:seed_variance}
\end{figure}
